\documentclass[11pt]{article}
\usepackage[T1]{fontenc}
\usepackage{lmodern}
\usepackage{amsmath,amssymb,amsthm,mathtools}
\usepackage[margin=1in]{geometry}
\usepackage{microtype}
\usepackage{booktabs,tabularx}
\usepackage{xcolor}
\usepackage[colorlinks=true,linkcolor=blue!45!black,citecolor=blue!45!black,urlcolor=blue!45!black]{hyperref}
\hypersetup{pdftitle={An explicit 2^3200 commutator bound via simultaneous block elimination},pdfauthor={Proof revision}}
\newtheorem{theorem}{Theorem}[section]
\newtheorem{proposition}[theorem]{Proposition}
\newtheorem{lemma}[theorem]{Lemma}
\newtheorem{corollary}[theorem]{Corollary}
\theoremstyle{remark}
\newtheorem{remark}[theorem]{Remark}
\newcommand{\C}{\mathbb C}
\newcommand{\R}{\mathbb R}
\newcommand{\tr}{\operatorname{tr}}
\newcommand{\Rea}{\operatorname{Re}}
\newcommand{\Ima}{\operatorname{Im}}
\newcommand{\diag}{\operatorname{diag}}
\newcommand{\rank}{\operatorname{rank}}
\newcommand{\codim}{\operatorname{codim}}
\newcommand{\Span}{\operatorname{span}}
\newcommand{\norm}[1]{\left\lVert#1\right\rVert}
\newcommand{\comm}[2]{[#1,#2]}
\newcommand{\HM}{\mathrm{HM}}
\numberwithin{equation}{section}
\setcounter{tocdepth}{1}
\allowdisplaybreaks[1]
\title{An explicit $2^{3200}$ commutator bound\\via simultaneous block elimination}
\author{Informal proof with a Lean correspondence appendix}
\date{September 17, 2026\\\small Revision following the supplied referee report}

\begin{document}
\maketitle
\begin{abstract}
Every trace-zero complex matrix $A$ admits a representation $A=BC-CB$ in its
original dimension with $\norm B\norm C\le 2^{3200}\norm A$, where all norms
are Euclidean operator norms. Following the supplied referee report, two
simultaneous square-zero shears replace the successive outside-block eliminations.
This removes the iterated polynomial from the high trace-mass branch. The
trace-preserving low-mass paving and the same-dimension strong induction are
retained. We give the full informal commutator proof, the numerical estimates
leading to $2^{3200}$, and a correspondence to the completed Lean development.
\end{abstract}

\noindent\textbf{Scope of this revision.}
This document revises the commutator proof in the repository's
\texttt{paper/main.tex}, using Comment 3 of the supplied report dated
15 September 2026. The two-shear improvement is credited to that report. The other mathematical ingredients are retained from the existing
manuscript and its formalization. The separation-family theorem and the
report's optional corollaries are outside this document's scope.

\medskip
\noindent\textbf{Formal result.}
The new declaration
\texttt{NoEpsilon.commutator\_bound\_two\_pow\_3200}
has the explicit constant in its conclusion and no additional branch hypotheses.
Its factors act on the same $\C^n$ as $A$; Appendix~\ref{sec:formal} records
norm semantics, dependencies, and reproducible checks.

\tableofcontents
\clearpage

\section{Introduction and main theorem}
For $A\in M_n(\C)$, write $\norm A$ for its operator norm on Euclidean space and
$\comm BC=BC-CB$. The quantitative commutator problem asks how economically one
can represent a trace-zero matrix by a single commutator. Johnson, Ozawa, and
Schechtman proved that, for each $\varepsilon>0$, a bound
$K_\varepsilon n^\varepsilon\norm A$ is possible; their first factor can be chosen
normal \cite{JOS}. Ravichandran and Srivastava obtained the bound
$300\exp(9\sqrt{\log n})\norm A$ \cite[Corollary 2]{RS}. The theorem considered here
allows both factors to be arbitrary complex matrices.

\begin{theorem}[Main theorem]\label{thm:main}
With $K=2^{3200}$, for every integer $n\ge1$
and every $A\in M_n(\C)$ with $\tr A=0$, there exist $B,C\in M_n(\C)$ satisfying
\begin{equation}\label{eq:main}
 A=\comm BC,\qquad \norm B\norm C\le K\norm A.
\end{equation}
The same $K$ works for every $n$.
\end{theorem}

The construction has two branches. Set
\begin{equation}\label{eq:globalparameters}
 s=8192,\qquad R=s^2=2^{26},\qquad b=8sR=2^{42},\qquad t_0=2/R.
\end{equation}
A matrix with large trace norm in every rotated Hermitian direction is handled
directly, with a constant $K_H(t_0)$. Otherwise it admits $R$ trace-zero
compressions of norm at most $319/R$ times its own norm. Reassembly costs a factor
$8s$ on the largest child budget, and
\[
 8s\frac{319}{R}=\frac{2552}{8192}<\frac12.
\]
This strict inequality is the quantitative reason the induction closes.

\subsection{Conventions and elementary norm estimates}
All spaces are finite-dimensional complex Hilbert spaces. A compression to a
subspace $V$ means $P_VA|_V$; no invariance of $V$ is implicit. Rectangular blocks
carry their operator norms. The trace is unnormalized, and
$|H|=(H^2)^{1/2}$ for Hermitian $H$. We use
\[
 \Rea A=(A+A^*)/2,\qquad \Ima A=(A-A^*)/(2i),
 \qquad \chi(S)=\norm S\norm{S^{-1}}.
\]
Unitary conjugation preserves the norm and trace. A similarity multiplies a
commutator product budget by at most $\chi(S)^2$.

We shall use a basic block estimate repeatedly. If an operator $T=(T_{ij})$ has
zero diagonal in an orthogonal decomposition with $k$ blocks and
$\norm{T_{ij}}\le c$ for $i\ne j$, then
\begin{equation}\label{eq:offnorm}
 \norm T\le c(k-1).
\end{equation}
Indeed, for $x=(x_j)$, the vector of component norms of $Tx$ is bounded
coordinatewise by $c(\mathbf1\mathbf1^\top-I)(\norm{x_j})_j$. The scalar matrix
has Euclidean operator norm $k-1$ when $k\ge2$; the case $k=1$ is immediate.
Also, the norm of a block diagonal operator is the largest block norm, and the
norm of a general block matrix is at most the sum of its block norms.

\subsection{External inputs}
Besides finite-dimensional spectral theory, the Toeplitz--Hausdorff theorem,
compact convexity, and the Banach fixed-point theorem, the proof uses three
specific published inputs: the two-dimensional Steinitz rearrangement theorem,
the simultaneous hollowisation theorem of Damm and Fa\ss bender, and the
independent-vector theorem of Marcus, Spielman, and Srivastava. Their exact
forms are stated in Section~\ref{sec:low}.
The two-commutator decomposition in the next section is elementary and allows
its first factors to depend on the input matrix.

\section{An elementary two-commutator decomposition}\label{sec:sumtwo}
\begin{lemma}\label{lem:hermitian}
Let $m\ge1$. If $G\in M_m(\C)$ is Hermitian and $\tr G=0$, there are a unitary $u$ and a matrix
$V$ such that $G=\comm uV$ and $\norm V\le\norm G$.
\end{lemma}
\begin{proof}
For $m=1$, take $u=I$ and $V=0$. Otherwise diagonalize $G$ unitarily and denote
its real eigenvalues by $\lambda_1,\ldots,\lambda_m$. They sum to zero and lie in
$[-a,a]$, where $a=\norm G$. They can be ordered so that every partial sum
$s_j=\sum_{i=1}^j\lambda_i$ belongs to $[-a,a]$. To see this, if the current sum
is positive, choose a remaining negative term; such a term exists because the
remaining terms have negative sum. If the current sum is negative, choose a
positive term. At zero choose any remaining term. In each case the new sum
stays in $[-a,a]$. This constructs the required ordering, including zero terms.

In this ordered eigenbasis, let
\[
 u=\sum_{j=1}^{m-1}E_{j,j+1}+E_{m,1},\qquad
 V=\sum_{j=1}^{m-1}s_jE_{j+1,j}.
\]
The first matrix is unitary, $\norm V=\max_{j<m}|s_j|\le a$, and multiplication
gives
\[
 uV=\diag(s_1,\ldots,s_{m-1},0),\qquad
 Vu=\diag(0,s_1,\ldots,s_{m-1}).
\]
Their difference is $\diag(\lambda_1,\ldots,\lambda_m)$, since $s_m=0$.
Conjugate both factors back to the original basis.
\end{proof}

\begin{corollary}\label{cor:sumtwo}
Let $m\ge1$. Every trace-zero $H\in M_m(\C)$ has a representation
\begin{equation}\label{eq:sumtwo}
 H=\comm uV+\comm kT,\qquad
 u,k\text{ unitary},\qquad
 \max(\norm V,\norm T)\le\norm H.
\end{equation}
\end{corollary}
\begin{proof}
Write $H=H_1+iH_2$ with $H_1,H_2$ Hermitian, trace zero, and of norm at most
$\norm H$. Apply Lemma~\ref{lem:hermitian} to obtain
$H_1=\comm uV$ and $H_2=\comm kW$, and set $T=iW$.
\end{proof}

\begin{remark}
The quantifiers in \eqref{eq:sumtwo} are $\forall H\,\exists u,k,V,T$. We do not
assert that one pair $u,k$ works uniformly for all $H$ in a fixed dimension, or
that the factors vary continuously with $H$. Neither property is needed below:
the four factors are chosen once before the Riccati equation is solved.
\end{remark}

\section{Solving an identity-corner core}\label{sec:riccati}
For $a\ge1$, define the following scalar quantities:
\begin{equation}\label{eq:P}
\begin{aligned}
 p&=2a,& d_0&=5a,& q&=5a+1,& d_1&=14a+1,\\
 L_0&=q^2+(2a+1)q+d_1,& L_1&=2q+2a+1,&
 \lambda&=2\max(L_0,L_1),\\[-2pt]
 P(a)&=(1+q)^4(\lambda+3)\bigl((\lambda+5)p+2\bigr).
\end{aligned}
\end{equation}
All of these are finite and nondecreasing in $a$.

\begin{proposition}\label{prop:core}
Let $m\ge1$, $a\ge1$, and
\[
 M=\begin{pmatrix}A&I_m\\F&D\end{pmatrix},\qquad
 \tr M=0,\qquad \norm M\le a.
\]
Then $M$ is a single commutator in $M_{2m}(\C)$ with factor norm product at most
$P(a)$.
\end{proposition}
\begin{proof}
Since $\tr(A+D)=0$, Corollary~\ref{cor:sumtwo} supplies fixed factors such that
\[
 A+D=\comm uV+\comm kT,\qquad
 \norm u=\norm k=1,\qquad \norm V,\norm T\le p.
\]
Set
\begin{equation}\label{eq:defects}
 U=\lambda I+u,\qquad D_0=\comm uV-A,\qquad
 D_1=\comm kV+D_0u-k+\comm u{Tk}-F.
\end{equation}
The block compression bounds imply $\norm A,\norm D,\norm F\le a$, whence
\[
 \norm{D_0}\le2p+a=d_0,\qquad
 \norm{D_1}\le2p+d_0+1+2p+a=d_1.
\]
On the closed ball $\mathcal B=\{Q:\norm{Q-D_0}\le1\}$ consider
\begin{equation}\label{eq:fixedpoint}
 \Psi(Q)=D_0+\lambda^{-1}
       \bigl(Q^2-DQ+Q(A-u)+D_1\bigr).
\end{equation}
Every $Q\in\mathcal B$ has norm at most $q$, so
\[
 \norm{\Psi(Q)-D_0}\le L_0/\lambda\le\tfrac12.
\]
For $Q,Q'\in\mathcal B$, the noncommutative identity
$Q^2-Q'^2=(Q-Q')Q+Q'(Q-Q')$ gives
\[
 \norm{\Psi(Q)-\Psi(Q')}\le
 \frac{2q+2a+1}{\lambda}\norm{Q-Q'}\le\tfrac12\norm{Q-Q'}.
\]
The ball is complete. The Banach fixed-point theorem gives an exact fixed point
$Q\in\mathcal B$.

Define
\begin{equation}\label{eq:companions}
 B_c=\begin{pmatrix}U&I\\k&0\end{pmatrix},\qquad
 C_c=\begin{pmatrix}
 V&T\\ A+Q-\comm uV+Tk&V+I-UT
 \end{pmatrix},\qquad
 S=\begin{pmatrix}I&0\\Q&I\end{pmatrix}.
\end{equation}
To verify the order of all factors, put
\[
 F_0=\comm kV+D_0U-k+\comm u{Tk}
     =\lambda D_0+D_1+F.
\]
Direct multiplication, using $A+D=\comm uV+\comm kT$, yields
\begin{equation}\label{eq:fourblocks}
\begin{aligned}
 \comm{B_c}{C_c}
  &=\begin{pmatrix}A+Q&I\\F_0-QU&D-Q\end{pmatrix},\\
 S^{-1}MS
  &=\begin{pmatrix}A+Q&I\\F+DQ-QA-Q^2&D-Q\end{pmatrix}.
\end{aligned}
\end{equation}
The bottom-left blocks agree precisely when
\[
 Q^2-DQ+Q(A-U)+F_0-F=0,
\]
which is equivalent to \eqref{eq:fixedpoint}. Consequently
\begin{equation}\label{eq:corecomm}
 M=\comm{SB_cS^{-1}}{SC_cS^{-1}}.
\end{equation}

It remains to bound the factors. The block norm sum gives
$\norm{B_c}\le\lambda+3$. The bottom-left block of $C_c$ equals $Q-D_0+Tk$,
of norm at most $1+p$, and the bottom-right block has norm at most
$1+(\lambda+2)p$. Thus
\[
 \norm{C_c}\le(\lambda+5)p+2,\qquad
 \norm S,\norm{S^{-1}}\le1+q.
\]
Equation~\eqref{eq:corecomm} now gives the claimed $P(a)$.
\end{proof}

\section{Absorbing a bounded number of remaining blocks}\label{sec:shears}
\begin{lemma}[Sylvester equation with separated centers]\label{lem:sylvester}
Let $U$ and $V$ act on possibly different Hilbert spaces, and let $z,w\in\C$.
If $\eta=\norm U+\norm V<|z-w|$, then, for every rectangular $Y$, the equation
\[
 (z-w)X+UX-XV=Y
\]
has a solution with
$\norm X\le\norm Y/(|z-w|-\eta)$.
\end{lemma}
\begin{proof}
On the space of rectangular operators with its operator norm, the map
$\mathcal E(X)=UX-XV$ has norm at most $\eta$. The inverse of
$(z-w)I+\mathcal E$ is the norm-convergent Neumann series
\[
 \frac1{z-w}\sum_{j=0}^\infty
       \left(-\frac{\mathcal E}{z-w}\right)^j,
\]
whose norm is at most $(|z-w|-\eta)^{-1}$.
\end{proof}

\begin{proposition}\label{prop:absorb}
Let a trace-zero matrix $A$ have an orthogonal block decomposition of sizes
$r,r,s_3,\ldots,s_b$, where $b\ge2$, $r\ge1$, and $1\le s_j\le r$.
Suppose $A_{12}=I_r$ and
$\norm A\le a$, with $a\ge1$. Put
\begin{equation}\label{eq:iterate}
\begin{gathered}
 q=b-2,\qquad u=1+a+\sqrt q,\qquad a_1=au^2,\\
 v=1+\sqrt q\,a_1,\qquad \widehat a=a_1v^2,\qquad k=q+1.
\end{gathered}
\end{equation}
Then $A$ has a same-dimensional single-commutator representation of cost at most
\begin{equation}\label{eq:absorb}
 (\widehat a/a)^2 k\bigl(4P(\widehat a)+2k\widehat a\bigr).
\end{equation}
\end{proposition}
\begin{proof}
Let $F=\bigoplus_{j=3}^b\C^{s_j}$. Choose isometries
$J_j:\C^{s_j}\to\C^r$ and form the row operator
$J=(J_3\ \cdots\ J_b):F\to\C^r$. Its summands need not have orthogonal
ranges. Nevertheless,
\[
 JJ^*=\sum_{j=3}^b J_jJ_j^*\le qI_r,\qquad \norm J\le\sqrt q.
\]
Use three large blocks of sizes $r,r,\dim F$, and put
\[
 X=J-A_{1F},\qquad
 S_1=\begin{pmatrix}I&0&0\\0&I&X\\0&0&I\end{pmatrix},\qquad
 A'=S_1^{-1}AS_1.
\]
The off-diagonal part of $S_1$ squares to zero, so changing the sign of $X$
gives its inverse. Left multiplication leaves the first block row fixed, and
right multiplication therefore gives
\[
 A'_{1F}=A_{1F}+A_{12}X=J,\qquad A'_{12}=I_r.
\]
In particular, $A'_{1j}=J_j$ for every outside block $j$.

Now form a column operator $Y:\C^r\to F$ with block rows
$Y_j=A'_{jj}J_j^*$, and set
\[
 S_2=\begin{pmatrix}I&0&0\\0&I&0\\Y&0&I\end{pmatrix},\qquad
 \widehat A=S_2^{-1}A'S_2.
\]
Again the inverse is obtained by changing the sign of the off-diagonal part.
Right multiplication changes only block column 1. Left multiplication
subtracts $Y_j$ times the first block row from each outside row $j$. Thus,
\begin{equation}\label{eq:shearzero}
 \widehat A_{jj}=A'_{jj}-Y_jA'_{1j}
 =A'_{jj}-A'_{jj}J_j^*J_j=0\quad(j\ge3),\qquad
 \widehat A_{12}=I_r.
\end{equation}
These identities account for interactions between distinct outside blocks:
no successive elimination or disjointness of the ranges of $J_j$ is assumed.

The rectangular compression satisfies $\norm{A_{1F}}\le a$, so
\[
 \norm X\le a+\sqrt q,\qquad
 \norm{S_1},\norm{S_1^{-1}}\le u,\qquad \norm{A'}\le a_1.
\]
For $z\in\C^r$ the orthogonal sum defining $F$ gives
\[
 \norm{Yz}^2=\sum_{j=3}^b\norm{A'_{jj}J_j^*z}^2
 \le q a_1^2\norm z^2.
\]
Hence $\norm Y\le\sqrt q\,a_1$ and
$\norm{S_2},\norm{S_2^{-1}}\le v$. Taking $S=S_1S_2$ yields
\begin{equation}\label{eq:telescope}
 \widehat A=S^{-1}AS,\qquad \norm{\widehat A}\le au^2v^2=\widehat a,
 \qquad \chi(S)\le u^2v^2=\widehat a/a.
\end{equation}
If $q=0$, both off-diagonal operators are zero and both shears are the
identity; the displayed upper bounds remain valid. Empty outside slots may
also be included, which is convenient for the formal fixed-slot partition.

Merge the first two blocks into one core $M$. Since every other diagonal block
is zero and similarity preserves trace, $\tr M=0$. Proposition~\ref{prop:core}
gives $M=\comm{B_0}{C_0}$ with product at most $P(\widehat a)$. Its identity
corner ensures $M\ne0$. Rescale the two factors inversely to arrange
\[
 \norm{B_0}=1/4,\qquad \norm{C_0}\le4P(\widehat a).
\]
There are now $k$ blocks. Let
\[
 \widehat B=\diag(B_0,I,2I,\ldots,(k-1)I),\qquad
 \widehat C_{11}=C_0,\quad \widehat C_{jj}=0\ (j>1).
\]
For distinct blocks solve
$\widehat B_{ii}\widehat C_{ij}-\widehat C_{ij}\widehat B_{jj}=\widehat A_{ij}$
using Lemma~\ref{lem:sylvester}. The centers are distinct integers and the only
non-scalar perturbation has norm $1/4$. Thus
$\norm{\widehat C_{ij}}\le2\widehat a$ and
\[
 \norm{\widehat B}\le k,\qquad
 \norm{\widehat C}\le4P(\widehat a)+2(k-1)\widehat a
\]
by \eqref{eq:offnorm}. These choices give
$\comm{\widehat B}{\widehat C}=\widehat A$. Conjugating back and applying
\eqref{eq:telescope} proves \eqref{eq:absorb}.
\end{proof}

\section{The high trace-mass branch}\label{sec:high}
For $0<t\le1$, say that $A\in M_n(\C)$ satisfies $\HM(t)$ if
\begin{equation}\label{eq:HM}
 \tr|\Rea(e^{i\theta}A)|\ge tn\norm A\quad\text{for every }\theta\in\R.
\end{equation}

\begin{lemma}[A disk in compressed numerical ranges]\label{lem:disk}
Suppose $\tr A=0$, $\norm A=1$, and $A$ satisfies $\HM(t)$. If $W$ has
codimension $d\le tn/4$, then the numerical range of $T=P_WA|_W$ contains the
closed disk of radius $\rho=t/4$ centered at zero.
\end{lemma}
\begin{proof}
For every $\theta$, the Hermitian matrix $H_\theta=\Rea(e^{i\theta}A)$ has norm
at most one and trace zero. Its positive eigenvalues therefore have sum at
least $tn/2$. Write the eigenvalues in decreasing order. Since $d<tn/2$, at
least $d+1$ eigenvalues are positive; moreover
\[
 tn/2\le d+(n-d)\lambda_{d+1}(H_\theta).
\]
It follows that
\[
 \lambda_{d+1}(H_\theta)\ge
 \frac{tn/2-d}{n-d}\ge t/4.
\]
The min--max principle yields
$\lambda_{\max}(P_WH_\theta|_W)\ge t/4$. The numerical range is compact and
convex by the Toeplitz--Hausdorff theorem. Its support in each direction is
this largest eigenvalue, so it contains the stated disk. Equivalently, use the
supporting-half-plane description of the ordinary numerical range; see
\cite[Section 2]{LiSze}.
\end{proof}

\begin{lemma}[A neutral vector with a large image]\label{lem:neutral}
If the numerical range of $T$ contains both $\rho$ and $-\rho$, where $\rho>0$,
then there is a unit vector $v$ such that
\begin{equation}\label{eq:neutral}
 v^*Tv=0,\qquad \norm{Tv}\ge\rho.
\end{equation}
\end{lemma}
\begin{proof}
Choose unit vectors $x_+,x_-$ with expectations $\rho,-\rho$. The density matrix
$D_*=(x_+x_+^*+x_-x_-^*)/2$ has trace one, satisfies $\tr(TD_*)=0$, and obeys
$\tr(T^*TD_*)\ge\rho^2$ by Cauchy--Schwarz.

On the compact convex set
\[
 \mathcal D=\{D\ge0:\tr D=1,\ \tr((\Rea T)D)=0,
                  \ \tr((\Ima T)D)=0\},
\]
maximize the real linear functional $\tr(T^*TD)$. The set of maximizers is a
nonempty compact face, so it has an extreme point that is also extreme in
$\mathcal D$. Every extreme point of $\mathcal D$ has rank one. In fact, if an
extreme $D$ had rank $s\ge2$, the Hermitian operators on its support would form
a real vector space of dimension $s^2\ge4$. The three displayed real linear
constraints admit a nonzero supported Hermitian perturbation $E$ in their
common kernel. Because $D$ is strictly positive on its support, both
$D+\eta E$ and $D-\eta E$ belong to $\mathcal D$ for sufficiently small
$\eta>0$, a contradiction. The maximizing extreme point is thus $vv^*$, and
its objective value is at least that of $D_*$.
\end{proof}

\begin{proposition}\label{prop:high}
Every trace-zero matrix satisfying $\HM(t)$ has a same-dimensional single
commutator representation with factor product at most $K_H(t)\norm A$, where
\begin{equation}\label{eq:KH}
\begin{gathered}
 L=\lceil16/t\rceil,\quad a_0=4/t,\quad u_*=1+a_0+\sqrt L,\\
 a_{1,*}=a_0u_*^2,\quad a_*=a_{1,*}(1+\sqrt L\,a_{1,*})^2,\quad k_*=L+1,\\
 K_H(t)=a_*^2k_*\bigl(4P(a_*)+2k_*a_*\bigr)<\infty.
\end{gathered}
\end{equation}
\end{proposition}
\begin{proof}
The zero matrix is immediate. Normalize $\norm A=1$, and set
$\rho=t/4$ and $r=\lceil tn/16\rceil$. We construct orthonormal
$v_1,\ldots,v_r$ such that all $v_i^*Av_j$ vanish and the vectors $Av_i$ are
pairwise orthogonal with norms at least $\rho$.

After selecting $v_1,\ldots,v_j$, work in
\[
 W_j=\Span\{v_i,Av_i,A^*v_i,A^*Av_i:1\le i\le j\}^{\perp}.
\]
For $0\le j\le r-1$ we have
$\codim W_j\le4j<tn/4$. Lemmas~\ref{lem:disk} and \ref{lem:neutral}, applied
to $P_{W_j}A|_{W_j}$, give a unit $v_{j+1}\in W_j$ with
\[
 v_{j+1}^*Av_{j+1}=0,\qquad
 \norm{Av_{j+1}}\ge\norm{P_{W_j}Av_{j+1}}\ge\rho.
\]
Orthogonality to $v_i$ preserves orthonormality; orthogonality to $Av_i$ and
$A^*v_i$ makes both mixed matrix coefficients vanish; orthogonality to
$A^*Av_i$ makes $Av_{j+1}$ perpendicular to $Av_i$. This proves the induction.

Let $V=\Span\{v_i\}$ and $W=A(V)$. The construction implies $V\perp W$ and
$\dim V=\dim W=r$; in particular $2r\le n$. In the ordered bases
$w_i=Av_i/\norm{Av_i}$ of $W$, followed by $v_i$ of $V$, the $(1,2)$ block of
$A$ is
\[
 D=\diag(\norm{Av_1},\ldots,\norm{Av_r}),\qquad \rho I\le D\le I.
\]
The similarity $S_0=\diag(I_r,D^{-1},I_{\rm rest})$ has condition number at
most $4/t$. In $S_0^{-1}AS_0$, the $(1,2)$ block equals $I_r$ and the norm is
at most $a_0=4/t$.

Partition the remaining space into blocks of positive size at most $r$. Since
$r\ge tn/16$, the number of outside blocks is at most $L=\lceil16/t\rceil$.
Apply Proposition~\ref{prop:absorb}. Every factor in its new budget is
nondecreasing in the nonnegative outside-block count, so the norm bound is at most $a_*$,
and after including $S_0$ the total similarity condition number is at most
$(4/t)(a_*/a_0)=a_*$. There are at most $k_*$ blocks after merging the core.
Its proof gives precisely the bound \eqref{eq:KH}. Finally undo the scalar
normalization.
\end{proof}

\section{The low trace-mass branch}\label{sec:low}
The key point of this branch is a complete paving whose individual traces are
exactly zero and whose ranks are equal.

\subsection{Rearrangement, hollowisation, and vector selection}
We use the following form of the Steinitz theorem: if $u_1,\ldots,u_N\in\R^2$
sum to zero and $\norm{u_i}_\infty\le1$, they can be reordered so that every
partial sum has $\ell_\infty$ norm at most $2$. This is the $d=2$ instance of
\cite[Theorem 1.2]{Barany}.

Damm and Fa\ss bender prove that three trace-zero Hermitian matrices can be
simultaneously conjugated by a unitary so that the first two have zero diagonal
and the third has at most two nonzero diagonal entries
\cite[Proposition 13(b)]{DF}. The last two entries need not vanish.

\begin{lemma}\label{lem:flatbasis}
Let $H,G$ be Hermitian operators on $\C^R$, with $R\ge2$, and let $E\ge0$.
There is an orthonormal basis $v_1,\ldots,v_R$ such that
\begin{equation}\label{eq:flatbasis}
 v_a^*Hv_a=\frac{\tr H}{R},\quad
 v_a^*Gv_a=\frac{\tr G}{R},\quad
 v_a^*Ev_a\le\frac{2\tr E}{R}\quad(1\le a\le R).
\end{equation}
\end{lemma}
\begin{proof}
Subtract the scalar mean from each of the three matrices and apply the cited
simultaneous hollowisation theorem. In the resulting basis the first two means
are attained at every vector. All but possibly two of the diagonal entries of
$E$ equal $\mu=\tr E/R$. The remaining two entries have sum $2\mu$ and are
nonnegative; each is therefore at most $2\mu$. This argument also covers
$R=2$.
\end{proof}

We also use the following theorem of Marcus, Spielman, and Srivastava
\cite[Theorem 1.4]{MSS}. If independent, finitely supported random vectors $x_i$
in a complex Hilbert space satisfy
\begin{equation}\label{eq:MSSassumptions}
 \sum_i\mathbb E x_ix_i^*=I,\qquad
 \mathbb E\norm{x_i}^2\le\varepsilon,
\end{equation}
where $\varepsilon>0$, some outcome satisfies
\begin{equation}\label{eq:MSSconclusion}
 \norm{\sum_i x_ix_i^*}\le(1+\sqrt\varepsilon)^2.
\end{equation}
The same conclusion holds if the covariance sum is merely at most $I$.
Indeed, spectrally decompose its positive semidefinite deficit and split each
eigenvalue into finitely many pieces at most $\varepsilon$. Add the corresponding
deterministic vectors, apply the theorem, and then discard their positive
semidefinite contributions. No upper restriction on $\varepsilon$ is required.

\subsection{Paving with exact quotas in every group}
\begin{lemma}\label{lem:transversals}
Let $D\ge0$ and let $h$ be a nonnegative integer. Suppose vectors are indexed
by finitely many groups of size $R=2^h$ and satisfy
\[
 \sum_i f_if_i^*\le I,\qquad \norm{f_i}^2\le D/R.
\]
Their indices can be partitioned into $R$ transversals, each containing exactly
one index from each group, such that every transversal $J$ satisfies
\begin{equation}\label{eq:transversalbound}
 \norm{\sum_{i\in J}f_if_i^*}\le\frac{C_D}{R},\qquad
 C_D=\bigl[1+2(\sqrt2+1)\sqrt D\bigr]^2.
\end{equation}
\end{lemma}
\begin{proof}
If $D=0$ all the vectors vanish, so any transversal partition works. Assume
$D>0$ and write $\delta=D/R$. Consider a node containing equally many indices
from each group, an even number per group, whose frame operator is at most
$LI$. Pair the indices within each group. For a pair of vectors $a,b$, choose
independently and uniformly one of the four vectors
\begin{equation}\label{eq:fourchoices}
 \sqrt2(a\oplus b),\quad \sqrt2(a\oplus(-b)),\quad
 \sqrt2(b\oplus a),\quad \sqrt2(b\oplus(-a)).
\end{equation}
For the random vector $w$ so obtained, the off-diagonal expected covariance
blocks cancel. On summing over the pairs,
\[
 \sum\mathbb E ww^*
   =\diag\left(\sum_i f_if_i^*,\sum_i f_if_i^*\right)\le LI,
 \qquad \mathbb E\norm w^2\le4\delta.
\]
When $L>0$, apply \eqref{eq:MSSconclusion} to $w/\sqrt L$ with
$\varepsilon=4\delta/L$. Some outcome has
\[
 \norm{\sum ww^*}\le(\sqrt L+2\sqrt\delta)^2.
\]
Its two diagonal blocks equal twice the frame operators of two complementary
halves, each containing one member of each pair. Consequently both children
have frame operator bounded by
\begin{equation}\label{eq:binaryrecurrence}
 L_{\rm next}I,\qquad L_{\rm next}=\tfrac12(\sqrt L+2\sqrt\delta)^2.
\end{equation}
If $L=0$, the vectors at the node are zero and can be divided arbitrarily.

Repeat this operation at every node for $h$ levels, starting with $L_0=1$.
At depth $j$ every node contains $R/2^j$ indices from each original group.
Thus the leaves are disjoint transversals covering all indices. Taking square
roots in \eqref{eq:binaryrecurrence} and summing the geometric progression gives
\[
 \sqrt{L_h}
 \le 2^{-h/2}+\sqrt{2\delta}\sum_{j=0}^{h-1}2^{-j/2}
 =R^{-1/2}+2(\sqrt2+1)\sqrt{D/R}(1-R^{-1/2}),
\]
which implies \eqref{eq:transversalbound}. The construction partitions indices,
so repeated or zero vector values cause no ambiguity.
\end{proof}

\subsection{Trace-preserving small compressions}
\begin{proposition}\label{prop:low}
Let $h\ge1$ and $k\ge1$ be integers, set $R=2^h$ and $n=kR$, and suppose
$\tr A=0$ and, for some
$\theta\in\R$,
\begin{equation}\label{eq:lowmass}
 \tr|\Rea(e^{i\theta}A)|\le2k\norm A.
\end{equation}
There are $R$ mutually orthogonal rank-$k$ projections $P_\ell$ summing to $I$
such that
\begin{equation}\label{eq:lowconclusion}
 \tr(P_\ell AP_\ell)=0,\qquad
 \norm{P_\ell AP_\ell}\le\frac{319}{R}\norm A.
\end{equation}
\end{proposition}
\begin{proof}
The case $A=0$ is immediate. Normalize and rotate so that
$\norm A=1$ and $A=H+iG$, where $H,G$ are trace-zero Hermitian matrices and
\[
 E=|H|\le I,\qquad \tr E\le2k.
\]
Choose an orthonormal eigenbasis $e_i$ of $G$, with eigenvalues $\lambda_i$.
Set
\[
 d_i=e_i^*Ee_i,\qquad \mu=\frac{\tr E}{kR}.
\]
Since $|\lambda_i|\le1$, $0\le d_i,\mu\le1$, the real vectors
$(\lambda_i,d_i-\mu)$ sum to zero and have $\ell_\infty$ norm at most one.
Apply Steinitz rearrangement and divide the reordered list into $k$ consecutive
groups of size $R$. Each group sum is a difference of two partial sums. The
corresponding orthogonal spectral subspaces $V_j$ therefore satisfy
\begin{equation}\label{eq:grouptraces}
 |\tr G_j|\le4,\qquad \tr E_j\le R\mu+4\le6,
\end{equation}
where subscripts denote compression to $V_j$. The spaces $V_j$ reduce $G$.

Apply Lemma~\ref{lem:flatbasis} to $H_j,G_j,E_j$ in each $V_j$. This gives an
orthonormal basis $v_{j,a}$ with
\begin{equation}\label{eq:groupbasis}
 v_{j,a}^*Hv_{j,a}=\tr H_j/R,\qquad
 v_{j,a}^*Gv_{j,a}=\tr G_j/R,\qquad
 v_{j,a}^*Ev_{j,a}\le12/R.
\end{equation}
For $f_{j,a}=E^{1/2}v_{j,a}$ we have
\[
 \sum_{j,a}f_{j,a}f_{j,a}^*=E\le I,\qquad \norm{f_{j,a}}^2\le12/R.
\]
Lemma~\ref{lem:transversals} with $D=12$ gives $R$ transversals. Let $P_\ell$
project onto the span of the original basis vectors in transversal $\ell$.
These projections are orthogonal, sum to $I$, and have rank $k$.

Equation~\eqref{eq:groupbasis} implies that the trace of each compression of
$A$ is
\[
 \frac1R\sum_j(\tr H_j+i\tr G_j)=\frac{\tr A}{R}=0.
\]
Since distinct $V_j$ reduce $G$, the compression $P_\ell GP_\ell$ is diagonal
in its selected basis, and \eqref{eq:grouptraces} gives
$\norm{P_\ell GP_\ell}\le4/R$. Furthermore
\[
 \norm{P_\ell EP_\ell}
 =\norm{E^{1/2}P_\ell E^{1/2}}
 =\norm{\sum_{(j,a)\in\ell}f_{j,a}f_{j,a}^*}\le C_{12}/R.
\]
The first equality follows from $\norm{X^*X}=\norm{XX^*}$ for
$X=E^{1/2}P_\ell$. Since $-E\le H\le E$, order and the Hermitian norm formula
imply $\norm{P_\ell HP_\ell}\le C_{12}/R$. Thus
\[
 \norm{P_\ell AP_\ell}\le(C_{12}+4)/R<319/R.
\]
For completeness, $\sqrt2<1.415$ and $\sqrt{12}<3.465$ give
$1+2(\sqrt2+1)\sqrt{12}<17.736$, whose square is less than $315$.
Undo the rotation and normalization to obtain \eqref{eq:lowconclusion}.
\end{proof}

\section{Finite-block assembly and the global induction}\label{sec:assembly}
\begin{lemma}\label{lem:assembly}
With $s,R,b$ as in \eqref{eq:globalparameters}, suppose $A$ is decomposed into
$1\le r\le2R-1$ nonempty orthogonal blocks and every diagonal block has a
representation $A_{ii}=\comm{U_i}{V_i}$ of cost
$p_i=\norm{U_i}\norm{V_i}$. Then $A=\comm BC$ in its original dimension, with
\begin{equation}\label{eq:assembly}
 \norm B\norm C\le8s\max_i p_i+b\norm A.
\end{equation}
\end{lemma}
\begin{proof}
For a nonzero diagonal block, inversely rescale its two factors to make
$\norm{U_i}=1/4$ and $\norm{V_i}=4p_i$. For a zero block use two zero factors;
this only decreases its cost. Choose distinct centers $z_i$ from the square grid
\[
 \{x+iy:x,y\in\{-s+\tfrac12,-s+\tfrac32,\ldots,s-\tfrac12\}\}.
\]
There are $4s^2=4R$ centers, with pairwise distance at least one. Define
\[
 B=\diag(z_iI+U_i),\qquad C_{ii}=V_i.
\]
Since $|z_i|\le\sqrt2(s-1/2)$, we have $\norm B<2s$. For each $i\ne j$ solve
\[
 (z_i-z_j)C_{ij}+U_iC_{ij}-C_{ij}U_j=A_{ij}.
\]
Lemma~\ref{lem:sylvester} applies with perturbation norm at most $1/2$, giving
$\norm{C_{ij}}\le2\norm{A_{ij}}\le2\norm A$. The block equations give
$\comm BC=A$, and \eqref{eq:offnorm} implies
\[
 \norm C\le4\max_i p_i+2(r-1)\norm A.
\]
Multiply by $2s$ and use $r-1\le2R$ to obtain \eqref{eq:assembly}.
No normality assumption on any $U_i$ is used.
\end{proof}

\begin{proof}[Proof of Theorem~\ref{thm:main}]
Let $K_H=K_H(t_0)$ be defined by \eqref{eq:KH}. We first prove the bound
with the internal budget
\begin{equation}\label{eq:K}
 K_0=\max\{8sK_H+b,\,2b\}
  =\max\{65536K_H+2^{42},\,2^{43}\}.
\end{equation}
This is a finite absolute constant. In this induction write $K=K_0$; the
quantitative calculation below shows $K_0\le2^{3200}$.
We argue by strong induction on $n$.
The zero matrix has the required representation, which also covers dimension
one. For a nonzero $A$, put $a=\norm A>0$.

First choose an orthonormal basis in which $A$ has zero diagonal. Here is a
short justification. The average of the diagonal expectations in any
orthonormal basis is zero. By convexity of the numerical range, zero is a
vector expectation. Choose a unit vector with this expectation; the
compression to its orthogonal complement still has trace zero. Repeat in the
smaller dimension. This also follows from \cite[Corollary 14]{DF}.

If $n<R$, use the $n$ zero singleton diagonal blocks in
Lemma~\ref{lem:assembly}; the cost is at most $ba\le Ka$.
If $n\ge R$, write
\[
 n=kR+q,\qquad k\ge1,\quad0\le q<R.
\]
Retain $N=kR$ vectors of the zero-diagonal basis and write $A_0$ for the
compression to their span. The other $q$ vectors are retained as zero
singleton diagonal blocks. Then $\tr A_0=0$ and $\norm{A_0}\le a$. All
off-diagonal interactions remain part of $A$ and will be handled by the
assembly lemma.

If $A_0=0$, every diagonal block in this decomposition is zero, and the
assembly lemma applies directly. Otherwise there are two cases.

If $A_0$ satisfies $\HM(t_0)$, Proposition~\ref{prop:high} gives a
representation of $A_0$ of cost at most $K_H\norm{A_0}\le K_Ha$. Assembling
this block with the $q$ zero singletons gives cost at most
$(8sK_H+b)a\le Ka$. The condition is tested with $A_0$'s own norm. When
$q=0$, this step is still a direct application of the high-mass proposition;
it does not appeal to induction in dimension $n$.

If $A_0$ fails $\HM(t_0)$, some $\theta$ satisfies
\[
 \tr|\Rea(e^{i\theta}A_0)|
 <t_0N\norm{A_0}=2k\norm{A_0}.
\]
Proposition~\ref{prop:low} supplies $R$ trace-zero blocks $T_i$ of dimension
$k<n$, with $\norm{T_i}\le(319/R)a$. Strong induction gives commutator
representations of these blocks of cost at most $(319K/R)a$. Together with
the $q$ zero singletons there are $R+q\le2R-1$ blocks. Applying
Lemma~\ref{lem:assembly} gives
\[
 \norm B\norm C
 \le\left(\frac{8s\cdot319}{R}K+b\right)a
 =\left(\frac{2552}{8192}K+b\right)a
 \le Ka,
\]
because $2552/8192<1/2$ and $K\ge2b$. This closes the induction.
\end{proof}

\subsection{Auditing the explicit constant}
At $t_0=2^{-25}$ the high-mass construction has
\[
 a_0=2^{27},\qquad L=2^{29},\qquad k_*=L+1\le2^{30},\qquad \sqrt L\le2^{15}.
\]
The two-shear budgets in \eqref{eq:KH} therefore satisfy
\begin{equation}\label{eq:dyadicbudgets}
 u_*\le2^{28},\qquad a_{1,*}\le2^{83},\qquad
 1+\sqrt L\,a_{1,*}\le2^{99},\qquad 1\le a_*\le2^{281}.
\end{equation}
For clarity we bound the core polynomial directly. Its scale, as defined in
Section~\ref{sec:riccati}, is
\[
 \sigma(x)=2\max\{(5x+1)^2+(2x+1)(5x+1)+14x+1,\,12x+3\}.
\]
For $x\ge1$, the first expression in the maximum is
$35x^2+31x+3\le69x^2$, and the second is at most $15x^2$.
Consequently $\sigma(x)\le138x^2$ and
\begin{align*}
 P(x)&=(5x+2)^4(\sigma(x)+3)((\sigma(x)+5)2x+2)\\
 &\le (7x)^4(141x^2)(288x^3)
 \le2^{29}x^9.
\end{align*}
Thus $P(a_*)\le2^{2558}$. The pullback through the initial normalization
and the two shears costs at most
$a_0^2(a_*/a_0)^2=a_*^2$. Substitution in \eqref{eq:KH} gives
\begin{align*}
 K_H(t_0)
 &\le (2^{281})^2\bigl(4\cdot2^{30}\cdot2^{2558}
                 +2(2^{30})^2\cdot2^{281}\bigr)\\
 &=2^{3152}+2^{904}\le2^{3153}.
\end{align*}
The high-case assembly budget is at most
$2^{16}2^{3153}+2^{42}<2^{3170}$.
In the low case the cost is
\[
 \left(\frac{319}{1024}K+2^{42}\right)\norm A\le K\norm A
 \quad\text{whenever }K\ge2^{43}.
\]
It follows that the internal budget $K_0$ in \eqref{eq:K} is at most
$2^{3200}$, proving the explicit value in Theorem~\ref{thm:main}.
All these comparisons are exact inequalities in the Lean proof; no
floating-point arithmetic is used.

\begin{remark}[Scope of the bound]
The exponent 3200 is a convenient generous choice, not an optimal exponent.
The improvement is that precisely two outside-block shears suffice, independent
of the number of outside blocks. The factors need not be normal. Auxiliary
doubled spaces occur only in the vector-selection argument; the final
commutator representation always acts on the original space.
\end{remark}

\appendix
\section{Lean correspondence and reproducibility}\label{sec:formal}
\subsection{The strengthened statement and its norm}
The public entry point \texttt{CommutatorTheorem.lean} imports the new module
\path{CommutatorTheorem/NoEpsilon/ExplicitBound.lean}. It activates
\texttt{Matrix.Norms.L2Operator}. The final declaration, with no hidden
section assumptions, states
\begin{quote}\small\ttfamily
commutator\_bound\_two\_pow\_3200\\
(n : Nat) (A : Matrix (Fin n) (Fin n) Complex)\\
(hA : Matrix.trace A = 0) :\\
\hspace*{1em}exists B C : Matrix (Fin n) (Fin n) Complex,\\
\hspace*{2em}A = B * C - C * B and\\
\hspace*{2em}norm B * norm C <= (2 : Real) \^{} 3200 * norm A
\end{quote}
Here the displayed ASCII notation is a readable transcription; the audit file
prints the elaborated Lean type. In particular the constant precedes all
choices of dimension and matrix, the trace is unnormalized, and both witnesses
have exactly the original matrix type. The theorem also covers $n=0$.

The declaration with suffix
\texttt{\_euclideanOperatorNorm} uses
\texttt{Matrix.toEuclideanCLM} in each norm. Its proof is the matrix-norm
declaration itself, with no norm-equivalence factor: the norms are definitionally
the same Euclidean operator norm. Thus this is not an entrywise or Frobenius
norm estimate.

\subsection{Mathematical correspondence}
Names use namespace \texttt{NoEpsilon}. All files are in the directory
\path{CommutatorTheorem/NoEpsilon/}.
\begin{center}\small
\begin{tabularx}{\textwidth}{@{}p{.30\textwidth}X@{}}
\toprule
Mathematical step & File and declaration \\
\midrule
Orthogonal sum estimate & \texttt{SimultaneousAbsorption.lean}:\newline
\texttt{Absorption.orthogonal\_sum\_norm\_sq} \\
Two simultaneous shears & \texttt{SimultaneousAbsorption.lean}:\newline
\texttt{Absorption.eliminate\_simultaneously} \\
Core and outside assembly & \texttt{AbsorptionCore.lean}:\newline
\texttt{Absorption.identity\_corner\_absorption} \\
Diagonal bridge normalization & \texttt{DiagonalAbsorption.lean}:\newline
\texttt{Absorption.diagonal\_corner\_absorption} \\
High-mass branch & \texttt{HighMassTheorem.lean}:\newline
\texttt{highMass\_bounded\_commutator} \\
Exact low-mass paving & \texttt{LowMassTheorem.lean}:\newline
\texttt{normalized\_low\_mass\_paving} \\
Same-dimension induction & \texttt{GlobalInduction.lean}:\newline
\texttt{boundedCommutator\_of\_lowMassBlocksAt} \\
Explicit numerical estimates & \texttt{ExplicitBound.lean}:\newline
\texttt{highMassNormBudget\_fixed\_le},\newline
\texttt{globalNormBudget\_le\_explicit} \\
Theorem~\ref{thm:main} & \texttt{ExplicitBound.lean}:\newline
\texttt{commutator\_bound\_two\_pow\_3200} \\
\bottomrule
\end{tabularx}
\end{center}
The sequential elimination lemmas remain available for compatibility, but the
refined main proof uses the simultaneous theorem. Its generic coordinate
proof uses orthogonal isometries $V_i$, the projection
$E=\sum_iV_iV_i^*$, and the row $L=\sum_iJ_iV_i^*$. It proves
$\norm E\le1$ and $\norm L\le\sqrt q$, then takes
$N=Q(L-P^*AE)$ and
$M=(\sum_iV_i(V_i^*A'V_iJ_i^*))P^*$.
These are precisely the two shears in Proposition~\ref{prop:absorb}.
It proves $N^2=M^2=0$, exact block identities, trace preservation,
and the condition-number and norm bounds, including empty outside slots.

The low-mass result retains the existing proof through Steinitz grouping,
simultaneous hollowization, and paired MSS selection. The MSS development
passes through stability of positive-semidefinite matrix pencils, the mixed
characteristic polynomial barrier bound, and interlacing selection. These
are proved inputs, rather than additional assumptions of the final theorem.
The final induction explicitly retains leftover zero singletons in the
original dimension and tests the normalized retained compression against
its own norm.

\subsection{Build and trust audit}
This revision uses Lean \texttt{v4.30.0-rc1} and the pinned Mathlib revision
in \texttt{lake-manifest.json}. The exact Mathlib commit and SHA-256 hashes
of the changed proof sources are recorded in
\texttt{verification/explicit\_3200/manifest.json}.
The base repository commit is
\begin{center}\small\ttfamily
8560e3605900427b4b8f615e58cd2d488262ad07.
\end{center}
The revision is a local working-tree change, not a newly published commit.
The recorded source hashes identify the checked refinement. From the
repository root, reproduce the checks with
\begin{quote}\small\ttfamily
lake build\\
lake env lean verification/explicit\_3200/Audit.lean
\end{quote}
The audit prints the final types and the transitive axiom sets for the
simultaneous elimination, the numerical high-case estimate, and both final
norm formulations. The final declarations depend only on
\begin{center}\ttfamily
[propext, Classical.choice, Quot.sound].
\end{center}
There is no \texttt{sorryAx}, custom mathematical axiom, or native decision
axiom in these dependencies. Full build and audit output accompanies the
source under \texttt{verification/explicit\_3200/}. A type audit is included
because an axiom list alone would not rule out hidden mathematical parameters.

\begin{thebibliography}{9}
\bibitem{Barany}
I.~B\'ar\'any,
\emph{A matrix version of the Steinitz lemma},
arXiv:2308.10102v2 (2024).
\url{https://arxiv.org/abs/2308.10102v2}.

\bibitem{DF}
T.~Damm and H.~Fa\ss bender,
\emph{Simultaneous hollowisation, joint numerical range, and stabilization by noise},
arXiv:1910.08813v2 (2020).
\url{https://arxiv.org/abs/1910.08813v2}.

\bibitem{JOS}
W.~B.~Johnson, N.~Ozawa, and G.~Schechtman,
\emph{A quantitative version of the commutator theorem for zero trace matrices},
arXiv:1202.0986 (2012).
\url{https://arxiv.org/abs/1202.0986}.

\bibitem{LiSze}
C.-K.~Li and N.-S.~Sze,
\emph{Canonical forms, higher rank numerical ranges, totally isotropic subspaces,
and matrix equations},
arXiv:0706.1536 (2007).
\url{https://arxiv.org/abs/0706.1536}.

\bibitem{MSS}
A.~W.~Marcus, D.~A.~Spielman, and N.~Srivastava,
\emph{Interlacing families II: Mixed characteristic polynomials and the
Kadison--Singer problem},
arXiv:1306.3969v4 (2014).
\url{https://arxiv.org/abs/1306.3969v4}.

\bibitem{RS}
M.~Ravichandran and N.~Srivastava,
\emph{Asymptotically optimal multi-paving},
arXiv:1706.03737v2 (2017).
\url{https://arxiv.org/abs/1706.03737v2}.
\end{thebibliography}
\end{document}
